\par\bigskip
\noindent\textbf{5. Subesquemas jacobianos formales}

\medskip
Sea $X_0$ un esquema de tipo finito sobre un cuerpo $k$, y sea
$(R,\mathfrak X)$ una deformación formal versal de $X_0$. El espacio
$\mathfrak X$ no es un $R$-esquema ordinario, sino un esquema formal
$R$-ádico. El objeto de esta sección es precisar el significado del
«lugar no liso de $\mathfrak X$ sobre $R$».

Recordemos primero los ideales de Fitting de un módulo de tipo finito.
Sea $A$ un anillo conmutativo y sea $E$ un $A$-módulo de tipo finito.
Elijamos una presentación
\[
K\longrightarrow F\longrightarrow E\longrightarrow0,
\]
donde $F$ es localmente libre de rango constante $n$. Para
$p=0,1,2,\ldots$, definamos
\[
I_A^p(E)=\operatorname{Im}\left(
 \bigwedge^{n-p}K\otimes_A\bigwedge^{n-p}F^*
 \longrightarrow A\right).
\]
Estos ideales son independientes de la presentación elegida. Son
invariantes del $A$-módulo $E$, llamados sus \emph{ideales de
Fitting}. Las propiedades siguientes se deducen directamente de la
definición.

\medskip
\noindent\textbf{5.1.}
\textbf{(a)} Para todo homomorfismo de anillos
$A\longrightarrow B$, la aplicación canónica
\[
B\otimes_A A/I_A^p(E)
\xrightarrow{\ \sim\ }
B/I_B^p(B\otimes_AE)
\]
es un isomorfismo para todo $p$. En particular,
\[
I_{S^{-1}A}^p(S^{-1}E)=S^{-1}I_A^p(E)
\]
para todo subconjunto multiplicativo $S\subset A$.

%%%%%%%%%%  scan idx 126  |  folio 115  %%%%%%%%%%

\smallskip
\noindent\textbf{(b)} Hay inclusiones
\[
I_A^0(E)\subset I_A^1(E)\subset I_A^2(E)\subset\cdots.
\]
Para todo $x\in\operatorname{Spec}(A)$,
\[
I_A^p(E)_x=A_x
\qquad\text{para todo }p\geq\dim_{\kappa(x)}E(x),
\]
donde $E(x)=\kappa(x)\otimes_AE$.

\smallskip
\noindent\textbf{(c)} Un punto
$x\in\operatorname{Spec}(A)-\operatorname{Supp}(A/I_A^p(E))$ si y
solo si, para cualquier presentación (equivalentemente, para alguna)
\[
K\longrightarrow F\longrightarrow E\longrightarrow0,
\qquad F\simeq A^n,
\]
el módulo $K_x$ contiene una copia de $A_x^{n-p}$ que es un sumando
directo de $F_x$. En particular,
\[
\operatorname{Supp}(A/I_A^0(E))=\operatorname{Supp}(E).
\]

\smallskip
\noindent\textbf{(d)} Para cualesquiera dos $A$-módulos de tipo
finito,
\[
I_A^r(E_1\oplus E_2)
=\sum_{p+q=r}I_A^p(E_1)I_A^q(E_2).
\]

\smallskip
\noindent\textbf{(e)} Si $A$ es un anillo ádico noetheriano,
entonces
\[
A/I_A^p(E)=\varprojlim_\nu A_\nu/I_{A_\nu}^p(E_\nu),
\]
donde $A=\varprojlim_\nu A_\nu$ y $E_\nu=E\otimes_AA_\nu$.

La propiedad 5.1(a) globaliza la definición. Si $X$ es un esquema y
$\underline E$ es un $\mathcal O_X$-módulo cuasicoherente de
presentación finita, existen haces cuasicoherentes de ideales
$\underline I_X^p(\underline E)$, $p\geq0$, caracterizados en todo
abierto afín $U\subset X$ por
\[
\Gamma(U,\underline I_X^p(\underline E))
\xrightarrow{\ \sim\ }
I_{\Gamma(U,\mathcal O_X)}^p(\Gamma(U,\underline E)).
\]

Para un morfismo de esquemas $X\longrightarrow Y$, el módulo relativo
de diferenciales de K\"ahler
$\Omega_{X|Y}=\Omega^1_{X|Y}$ es cuasicoherente. Si el morfismo es
esencialmente de presentación finita, entonces $\Omega_{X|Y}$ es de
presentación finita, de modo que están definidos los haces de ideales
$\underline I_X^p(\Omega_{X|Y})$.

\medskip
\noindent\textbf{Definición 5.2.}
Sea $X\longrightarrow Y$ un morfismo esencialmente de presentación
finita. Para todo entero $p\geq0$, definamos $J_X^p(X|Y)$ como el
subesquema cerrado de $X$

%%%%%%%%%%  scan idx 127  |  folio 116  %%%%%%%%%%

\noindent definido por el haz de ideales
\[
\underline I_X^p(X|Y)=\underline I_X^p(\Omega_{X|Y}).
\]
Se llama el $p$-ésimo \emph{subesquema jacobiano} de $X$ sobre $Y$.
Para todo $x\in X$, pongamos
\[
X_x=\operatorname{Spec}(\mathcal O_{X,x}).
\]
